% =====================================================================
%  The Onix Protocol — single-source LaTeX manuscript (Ledger template)
%  RUSSIAN TRANSLATION of onix.tex (перевод английского источника).
%  Author: Anatoly Piskunov (ORCID 0009-0000-8883-4111)
%
%  TARGET CLASS: ledger.cls  (LedgerJournal/Ledger-Latex-Template, GitHub)
%    Requires ledger.cls, ledgerbib.bst and an images/ folder on the build
%    path.  Compile with pdflatex.  This same .tex is the arXiv source (arXiv
%    primary cs.GT, cross-list q-fin.TR).
%
%  NOTES FOR THE FINAL BUILD
%    * ledger.cls already loads amsmath, amssymb, amsthm, hyperref, cleveref,
%      enumitem, tabularx, titling, titlesec, mathptmx, microtype.  The
%      preamble below therefore stays minimal to avoid clashes.
%    * Русская сборка требует кириллических шрифтов: подключены пакеты
%      fontenc[T2A] / inputenc[utf8] / babel[russian,english].  Ledger
%      использует mathptmx; если Overleaf ругается на шрифты, автор может
%      подстроить fontenc/шрифтовую схему (например, взять pscyr или
%      lmodern) без изменения содержимого.
%    * References use a numbered thebibliography with parenthetical author-year
%      cites in prose (as in the source draft).  Ledger house style is endnotes
%      (\endnote/\theendnotes); converting is mechanical if the editor asks.
%    * No figures/images are used (the paper is math + tables only).
%
%  CHANGE OVER THE JULY 2026 DRAFT: this version weaves in the finalized F1
%  early-exit "deferred outcome-contingent claim" mechanism (new Section on
%  Early Exit + Theorem 2), reconciled with the leverage subsystem, governance
%  parameters, limitations and experimental hypotheses.  The mechanism is
%  live-verified on the VIZ testnet (full cycle: cancel -> claim recorded ->
%  settle paid at the capped bucket with a visible haircut -> claim consumed).
% =====================================================================
\documentclass{ledger}

% --- Russian language support (added for the ru build) ---
\usepackage[T2A]{fontenc}
\usepackage[utf8]{inputenc}

% --- theorem environments (amsthm is loaded by the class) ---
\theoremstyle{plain}
\newtheorem{theorem}{Теорема}
\newtheorem{corollary}{Следствие}
\theoremstyle{remark}
\newtheorem*{remark}{Замечание}

% --- Ledger front-matter counters / metadata (edit at typesetting) ---
\newcommand{\thefirstpagenum}{1}
\newcommand{\thelastpagenum}{1}
\newcommand{\theyear}{2026}
\newcommand{\thevol}{X}
\newcommand{\thedoi}{DOI 10.5195/LEDGER.\theyear.XXX}

\hypersetup{
  pdfauthor={Anatoly Piskunov},
  pdftitle={Объединение тотализаторного расчёта с автоматическими маркет-мейкерами для рынков предсказаний с гарантией ликвидности: протокол Onix}
}

% Clean the template banner: show only the title (remove "LEDGER LATEX TEMPLATE").
\renewcommand{\pretitle}{\begin{center}\fontsize{18pt}{22pt}\bfseries\selectfont}
\renewcommand{\posttitle}{\par\end{center}\vspace{0.6em}}

\title{Объединение тотализаторного расчёта с автоматическими маркет-мейкерами
для рынков предсказаний с гарантией ликвидности: протокол Onix}

\author{Анатолий Пискунов\thanks{Независимый исследователь --- распределённый
реестр VIZ. Email: \href{mailto:anatoly.piskunov@gmail.com}{anatoly.piskunov@gmail.com}.
ORCID: \href{https://orcid.org/0009-0000-8883-4111}{0009-0000-8883-4111}.}}

\date{2026}

\begin{document}
\selectlanguage{russian}
\maketitle

\begin{abstract}
Рынки предсказаний агрегируют рассеянную информацию в вероятностные прогнозы,
которые устойчиво превосходят опросы и экспертные панели, однако их
распространение сдерживается структурной проблемой ликвидности: поставщики
рыночной глубины несут риск неблагоприятного отбора, отпугивающий розничных
участников. Мы представляем \textbf{протокол Onix} --- гибридную архитектуру,
которая разделяет функцию \emph{ценообразования} и функцию \emph{расчёта} на
рынках предсказаний. Сочетая ценообразование автоматического маркет-мейкера ---
маркет-мейкера с постоянным произведением (CPMM) для бинарных исходов и
логарифмического правила рыночного скоринга (LMSR) для многоисходных рынков ---
с тотализаторным (parimutuel) расчётом, мы достигаем \textbf{структурной
гарантии} того, что основной капитал поставщика ликвидности никогда не
подвергается риску из-за исходов ставок. Мы формализуем экономические инварианты
протокола, доказываем гарантию основного капитала LP для обоих типов рынков и
разрешаем центральное напряжение любого рынка с ценой по кривой --- то, что
ранний выход через кривую ценообразования реализует торговую прибыль/убыток
(P\&L) против тех самых LP, которых протокол защищает, --- с помощью механизма
\emph{отложенного права требования, зависящего от исхода}, для которого доказываем
второй инвариант без эмиссии / безопасности LP. Мы описываем «ленивый» пул
ликвидности, обеспечивающий пассивное участие розничных пользователей,
анализируем механизм разрешения споров под управлением DAO и обсуждаем
экспериментальные гипотезы, для проверки которых спроектирована эта система.
Протокол реализован как операции консенсусного уровня в распределённом реестре
VIZ.
\end{abstract}

\begin{keywords}
\item{рынки предсказаний} \item{автоматические маркет-мейкеры} \item{тотализатор} \item{LMSR} \item{предоставление ликвидности} \item{управление через DAO} \item{дизайн механизмов}
\end{keywords}

\clearpage
\tableofcontents
\clearpage

% =====================================================================
\section{Введение}\label{sec:intro}

\subsection{Проблема ликвидности рынков предсказаний}

Рынки предсказаний --- одни из наиболее надёжных механизмов агрегирования
рассеянной информации в применимые вероятностные прогнозы. Эмпирические данные
из политического прогнозирования (Berg, Forsythe, Nelson \& Rietz, 2008),
спортивных ставок (Levitt, 2004) и внутрикорпоративных рынков (Cowgill, Wolfers
\& Zitzewitz, 2009) показывают, что рыночные цены устойчиво превосходят
альтернативные методы прогнозирования.

Несмотря на этот послужной список, развёртывание рынков предсказаний остаётся
ограниченным. Фундаментальное ограничение --- не спрос на прогнозы, а
\emph{предложение ликвидности}. Каждая существующая архитектура рынка
предсказаний требует, чтобы поставщики ликвидности (LP) приняли один из трёх
профилей риска:

\begin{enumerate}
  \item \textbf{Непостоянные потери} (impermanent loss) на рынках на основе AMM
  (Adams, Zinsmeister \& Robinson, 2020; Adams et al., 2021), где LP
  систематически проигрывают информированным трейдерам;
  \item \textbf{Неблагоприятный отбор} на рынках с книгой лимитных заявок, где
  маркет-мейкеры несут инвентарный риск против лучше информированных
  контрагентов;
  \item \textbf{Ограниченные потери субсидии} на рынках LMSR (Hanson, 2003), где
  маркет-мейкер рискует суммой до $b \ln N$ против корректно информированной
  «толпы».
\end{enumerate}

Эти профили риска ограничивают предоставление ликвидности кругом искушённых
участников со способностью к активному управлению, порождая тонкие рынки с
высоким проскальзыванием, которые отпугивают бетторов, --- самоподдерживающийся
цикл:

\[
\text{тонкие рынки} \rightarrow \text{высокое проскальзывание} \rightarrow \text{плохой UX}
\rightarrow \text{мало бетторов} \rightarrow \text{низкие комиссии} \rightarrow
\text{нет стимула для LP} \rightarrow \text{тонкие рынки}
\]

\subsection{Исследовательские вопросы}

Эта работа ставит и стремится ответить на следующие вопросы через развёрнутую
экспериментальную систему:

\textbf{Q1.} \emph{Можно ли архитектурно разделить функции ценообразования и
расчёта рынка предсказаний так, чтобы основной капитал LP был структурно
изолирован от исходов ставок?}

\textbf{Q2.} \emph{Устраняет ли тотализаторный механизм расчёта, где проигравшие
исключительно финансируют выигравших, проблему неблагоприятного отбора, присущую
рынкам предсказаний на основе AMM и LMSR?}

\textbf{Q3.} \emph{Способно ли автоматическое размещение ликвидности через
пуловый механизм («ленивый» пул) снизить барьер входа для LP настолько, чтобы
запустить рыночную глубину без активного выбора рынков отдельными LP?}

\textbf{Q4.} \emph{Является ли механизм разрешения споров под управлением DAO, со
взвешенным по стейку публичным голосованием и оракулами с залогом, жизнеспособной
альтернативой централизованному арбитражу или оптимистичным оракулам?}

\textbf{Q5.} \emph{Какая полезность токена управления возникает в протоколе рынка
предсказаний, где позиции LP несут вес голоса, и создаёт ли это устойчивую петлю
обратной связи спроса на токен?}

\subsection{Вклад}

Наш вклад состоит в следующем:

\begin{enumerate}
  \item \textbf{Формальное доказательство} того, что сочетание ценообразования
  AMM с тотализаторным расчётом даёт структурную гарантию основного капитала LP
  для обоих типов рынков --- бинарного (CPMM) и многоисходного (LMSR)
  (Теорема~\ref{thm:lp}).
  \item \textbf{Механизм отложенного права требования, зависящего от исхода},
  который разрешает напряжение «выход по кривой / тотализатор» --- позволяя
  бетторам и позициям с плечом выходить рано, не реализуя торговую P\&L против
  глубины LP, --- вместе с доказательством (Теорема~\ref{thm:earlyexit}) того, что
  он не может создать непокрытую недостачу, эмитировать токены или затронуть
  основной капитал LP.
  \item \textbf{Механизм «ленивого» пула} --- система автоматического размещения
  капитала с градуированным отзывом и защитой от издержек упущенных возможностей,
  обеспечивающая пассивное участие розничных LP, плюс опциональная, выключенная по
  умолчанию подсистема кредитного плеча, финансируемая из пула, чья экспозиция
  \emph{ограничена} и изолирована от гарантии LP по Теореме~\ref{thm:lp}.
  \item \textbf{Двухрежимная система разрешения споров}, сочетающая оракулов с
  залогом с арбитражем комитета DAO, включая полный теоретико-игровой анализ
  совместимости стимулов оракула.
  \item \textbf{Экспериментальный протокол}, реализованный как операции
  консенсусного уровня в продакшен-реестре распределённого реестра, с явными
  гипотезами и измеримыми результатами.
\end{enumerate}

% =====================================================================
\section{Связанные работы}\label{sec:related}

\subsection{Правила рыночного скоринга}

Hanson (2003) ввёл логарифмическое правило рыночного скоринга (LMSR), в котором
маркет-мейкер котирует цены через функцию softmax над параметрами количеств. LMSR
гарантирует ограниченные потери маркет-мейкера в размере $b \ln N$ и обеспечивает
сумму цен, равную единице. Однако эта граница одновременно является и
максимальным риском LP: если «толпа» верно предсказывает исход, маркет-мейкер
теряет всю субсидию. Это ограничило применение LMSR корпоративными средами
(Microsoft, Inkling), где оператор поглощает потери.

\subsection{Маркет-мейкеры с постоянным произведением}

Формула постоянного произведения $x \cdot y = k$, популяризированная Uniswap
(Adams, Zinsmeister \& Robinson, 2020), обеспечивает непрерывное ценообразование
для пар из двух активов. В применении к бинарным рынкам предсказаний резервы
представляют противоположные исходы, и инвариант даёт естественное вероятностное
ценообразование. Однако стандартные LP в AMM несут непостоянные потери ---
эмпирический анализ показывает, что более 50\% LP в Uniswap v3 уступают стратегии
«купи и держи» (Adams et al., 2021).

\subsection{Фреймворки условных токенов}

Polymarket использует фреймворк условных токенов Gnosis (CTF), где позиции --- это
токены ERC-1155 с операциями разделения/слияния для обеспечения согласованности
цен ($\sum p_i = 1$). Это добавляет компонуемость, но вносит сложность и требует
внешнего арбитражного механизма. Оптимистичный оракул UMA обеспечивает разрешение
споров через игру «вызов--ответ» с экономическими залогами.

\subsection{Тотализаторные ставки}

Тотализаторная (parimutuel) модель, восходящая к тотализатору Пьера Олле (Pierre
Oller, 1865), объединяет все ставки и распределяет ставки проигравших
пропорционально между выигравшими. Это гарантирует, что «дом» никогда не платит
из собственного капитала. Однако традиционные тотализаторные системы лишены
ценообразования в реальном времени --- коэффициенты окончательны лишь при закрытии
пула, --- что делает их непригодными для непрерывного агрегирования информации.

\subsection{Наш подход}

Протокол Onix сочетает способность AMM/LMSR к ценообразованию в реальном времени
со свойством нулевого риска LP, присущим тотализаторному расчёту. AMM или LMSR
служит исключительно \emph{механизмом ценообразования} --- определяет
подразумеваемые вероятности и назначает веса, --- тогда как расчёт следует
тотализаторной модели, где проигравшие финансируют выигравших, а основной капитал
LP архитектурно отделён от потока выплат.

Это позиционирует Onix иначе, чем два других гибридных направления в литературе.
Во-первых, \emph{чувствительные к ликвидности} маркет-мейкеры (Othman, Pennock,
Reeves \& Sandholm, 2013; Abernethy, Chen \& Vaughan, 2011) сохраняют оператора в
роли контрагента, но ограничивают или формируют его потери через геометрию функции
издержек; оператор всё ещё может проиграть. Onix же полностью убирает роль
контрагента у LP --- кривая лишь \emph{ценообразует}, она никогда не \emph{платит}.
Во-вторых, гибриды \emph{CLOB+AMM} (например, книги заявок, подстрахованные AMM,
как в ряде дизайнов DEX и в CLOB Polymarket с AMM-резервом) смешивают две площадки
ценообразования, но наследуют инвентарный риск и неблагоприятный отбор со стороны
маркет-мейкинга. Onix сохраняет единую непрерывную площадку ценообразования и
полностью выносит несение риска из ценообразования в тотализаторный пул. Насколько
нам известно, сочетание непрерывного ценообразователя AMM/LMSR с тотализаторным
расчётом (проигравшие финансируют выигравших) для получения безусловной гарантии
основного капитала LP ранее не формализовалось.

% =====================================================================
\section{Модель протокола}\label{sec:model}

\subsection{Участники и роли}

Мы определяем следующее множество участников $\mathcal{P}$:

\begin{center}
\begin{tabular}{lll}
\hline
\textbf{Роль} & \textbf{Символ} & \textbf{Функция} \\
\hline
Создатель рынка & $c \in \mathcal{P}$ & Задаёт параметры рынка, предоставляет начальную ликвидность \\
Оракул & $o \in \mathcal{P}$ & Вносит страховой залог $I_o$, принимает рынки, разрешает исходы \\
Беттор & $b_i \in \mathcal{P}$ & Делает ставки на исходы, получает токены исхода \\
Поставщик ликвидности & $\ell_j \in \mathcal{P}$ & Вносит капитал в пулы рынков или в «ленивый» пул \\
Арбитр спора & $\mathcal{D}$ & Комитет (взвешенное по стейку голосование) или именованный аккаунт \\
\hline
\end{tabular}
\end{center}

\subsection{Определение рынка}

Рынок $\mathcal{M}$ --- это кортеж:

\[
\mathcal{M} = (o, \tau, \mathbf{O}, t_{\text{bet}}, t_{\text{res}},
\boldsymbol{\theta}, L)
\]

где $o$ --- назначенный оракул, $\tau \in \{0, 1\}$ --- тип рынка (бинарный или
многоисходный), $\mathbf{O} = \{O_1, \ldots, O_N\}$ --- множество исходов
($N = 2$ для бинарного, $3 \leq N \leq 10$ для многоисходного), $t_{\text{bet}}$ и
$t_{\text{res}}$ --- времена истечения приёма ставок и разрешения,
$\boldsymbol{\theta}$ --- вектор параметров комиссий, $L$ --- начальная ликвидность.

\subsection{Структура комиссий}\label{sec:fees}

Все комиссии выражены в базисных пунктах (bp), где $10000 \text{ bp} = 100\%$.
Вектор комиссий:

\[
\boldsymbol{\theta} = (\theta_{\text{oracle}}, \theta_{\text{creator}},
\theta_{\text{liq}})
\]

Комиссии вычисляются исключительно при разрешении из совокупной ставки проигравшей
стороны $S_{\text{lose}}$ и никогда не удерживаются в момент размещения ставки:

\begin{align*}
f_{\text{oracle}}  &= \lfloor S_{\text{lose}} \cdot \theta_{\text{oracle}} / 10000 \rfloor \\
f_{\text{creator}} &= \lfloor S_{\text{lose}} \cdot \theta_{\text{creator}} / 10000 \rfloor \\
f_{\text{liq}}     &= \lfloor S_{\text{lose}} \cdot \theta_{\text{liq}} / 10000 \rfloor \\
W &= S_{\text{lose}} - f_{\text{oracle}} - f_{\text{creator}} - f_{\text{liq}}
\end{align*}

где $W$ --- \textbf{пул победителей}, доступный для распределения выигравшим
бетторам.

\subsection{Конечный автомат рынка}\label{sec:fsm}

Жизненный цикл $\mathcal{M}$ следует конечному автомату:

\[
q_0 \xrightarrow{\text{оракул принимает}} q_1 \xrightarrow{t \geq t_{\text{bet}}}
q_2 \xrightarrow{\text{оракул разрешает}} q_3 \xrightarrow{\text{льготный период}}
\text{выплачено}
\]

Из $q_0$ ведут два пути в терминальное удалённое состояние $q_{-1}$, оба возвращают
начальную ликвидность создателю (антиспам-комиссия за создание, уже уплаченная в
фонд DAO при создании, \emph{не} возвращается): явное отклонение
$q_0 \xrightarrow{\text{оракул отклоняет}} q_{-1}$ и \emph{таймаут}
$q_0 \xrightarrow{t \geq t_0 + \tau_{\text{accept}}} q_{-1}$, срабатывающий
поблочным кроном, когда оракул не принимает и не отклоняет рынок в течение
заданного управлением окна акцепта $\tau_{\text{accept}}$ (по умолчанию один час).
Таймаут ограничивает, как долго неотвечающий оракул может удерживать в заложниках
начальный капитал создателя.

% =====================================================================
\section{Механизмы ценообразования}\label{sec:pricing}

\subsection{Бинарные рынки: маркет-мейкер с постоянным произведением}

Для бинарных рынков ($N = 2$) протокол поддерживает два резерва $R_A$ и $R_B$,
подчинённых инварианту:

\[
k = R_A \cdot R_B
\]

\textbf{Инициализация.} При начальной ликвидности $L$:

\[
R_A = \lfloor L / 2 \rfloor, \quad R_B = L - R_A, \quad k = R_A \cdot R_B
\]

\textbf{Размещение ставки.} Когда беттор $b_i$ ставит сумму $a$ на исход $A$:

\begin{align*}
R'_A &= R_A + a \\
R'_B &= \lfloor k / R'_A \rfloor \\
w_i  &= R_B - R'_B \quad \text{(полученные токены)}
\end{align*}

Здесь $w_i$ --- \emph{вес} беттора. Это \textbf{вычисляемая} величина --- количество
токенов исхода, созданных кривой, --- и она в общем случае \textbf{не} равна
внесённой сумме $a$: в зависимости от текущих резервов $w_i$ может быть больше или
меньше $a$. Вес используется исключительно для распределения пула победителей при
расчёте (\Cref{sec:unified}); он никогда не является требованием, выраженным в
валюте.

\textbf{Подразумеваемая вероятность.} Подразумеваемая рынком вероятность каждого
исхода:

\[
P(A) = \frac{R_A}{R_A + R_B}, \quad P(B) = \frac{R_B}{R_A + R_B}
\]

Заметим, что $P(A) + P(B) = 1$ по построению, что устраняет необходимость во
внешнем арбитражном механизме для обеспечения согласованности цен.

\subsection{Многоисходные рынки: LMSR}

Для рынков с $N > 2$ исходами протокол использует логарифмическое правило
рыночного скоринга с ценообразованием softmax.

\textbf{Функция издержек:}

\[
C(\mathbf{q}) = b \cdot \ln\left(\sum_{j=1}^{N} \exp(q_j / b)\right)
\]

где $\mathbf{q} = (q_1, \ldots, q_N)$ --- параметры количеств, а $b$ --- параметр
ликвидности.

\textbf{Функция цены (softmax):}

\[
p_i(\mathbf{q}) = \frac{\exp(q_i / b)}{\sum_{j=1}^{N} \exp(q_j / b)}
\]

По определению softmax, $\sum_{i=1}^{N} p_i = 1$ тождественно.

\textbf{Издержки сделки.} Издержки покупки $\Delta$ токенов на исход $i$:

\[
\text{cost}(\Delta, i) = C(\mathbf{q} + \Delta \cdot \mathbf{e}_i) - C(\mathbf{q})
\]

\textbf{Параметр ликвидности.} Параметр $b$ финансируется субсидией LP $S$:

\[
b = \frac{S}{\ln N}
\]

\textbf{Численная устойчивость.} Вычисление log-sum-exp использует стандартное
тождество:

\[
\ln\left(\sum_j \exp(x_j)\right) = \max(\mathbf{x}) +
\ln\left(\sum_j \exp(x_j - \max(\mathbf{x}))\right)
\]

% =====================================================================
\section{Тотализаторный расчёт и гарантия LP}\label{sec:settlement}

\subsection{Единая модель расчёта}\label{sec:unified}

Оба типа рынков используют идентичный механизм расчёта. При разрешении оракул
объявляет выигравший исход $O^* \in \mathbf{O}$. Пусть $\mathcal{W}$ --- множество
ставок на $O^*$, а $\mathcal{L}$ --- все остальные ставки.

\[
S_{\text{lose}} = \sum_{b_i \in \mathcal{L}} a_i
\]

где $a_i$ --- сумма ставки. Пул победителей $W$ вычисляется как в \Cref{sec:fees}.

Каждый выигравший беттор $b_i \in \mathcal{W}$ получает:

\[
\pi_i = a_i + \frac{w_i}{\sum_{b_j \in \mathcal{W}} w_j} \cdot W - \tau_i
\]

где $w_i$ --- вес беттора (токены из CPMM или LMSR), а $\tau_i$ --- временной штраф
на прибыль (\Cref{sec:timepenalty}).

\subsection{Теорема: гарантия основного капитала LP}\label{sec:thm1}

\begin{theorem}[Гарантия основного капитала LP]\label{thm:lp}
При тотализаторном расчёте суммарная распределённая сумма равна суммарной
полученной сумме, причём основной капитал LP $L$ возвращается безусловно. Основной
капитал LP никогда не подвергается риску из-за исходов ставок.
\end{theorem}

\begin{proof}
Рассмотрим полный денежный поток через рынок:

\[
\text{Money}_{\text{IN}} = L + \sum_{b_i \in \mathcal{W}} a_i +
\sum_{b_i \in \mathcal{L}} a_i = L + \sum_{\text{all bets}} a_i
\]

где первая сумма $\sum_{b_i \in \mathcal{W}} a_i$ --- совокупный основной капитал,
поставленный \emph{выигравшими} (возвращается им полностью при расчёте), а
$S_{\text{lose}} = \sum_{b_i \in \mathcal{L}} a_i$ --- совокупная ставка, утраченная
\emph{проигравшими} (единственный источник пула победителей и всех комиссий). При
разрешении исходящие потоки таковы:

\[
\text{Money}_{\text{OUT}} = L + \sum_{b_i \in \mathcal{W}} a_i + W +
f_{\text{oracle}} + f_{\text{creator}} + f_{\text{liq}}
\]

Подставляя $W = S_{\text{lose}} - f_{\text{oracle}} - f_{\text{creator}} -
f_{\text{liq}}$:

\begin{align*}
\text{Money}_{\text{OUT}}
&= L + \sum_{b_i \in \mathcal{W}} a_i + S_{\text{lose}} \\
&= L + \sum_{b_i \in \mathcal{W}} a_i + \sum_{b_i \in \mathcal{L}} a_i
 = L + \sum_{\text{all bets}} a_i = \text{Money}_{\text{IN}} \qquad \qedhere
\end{align*}
\end{proof}

\begin{corollary}
Гарантия основного капитала LP выполняется независимо от механизма назначения
весов, числа бетторов, распределения ставок по исходам или параметров комиссий.
Это свойство архитектуры расчёта, а не кривой ценообразования.
\end{corollary}

\begin{corollary}
Максимальная суммарная выплата выигравшим равна $S_{\text{lose}}$ (весь пул
проигравших), независимо от весов, назначенных AMM или LMSR. Следовательно, выплаты
выигравшим никогда не могут черпаться из капитала LP.
\end{corollary}

\begin{remark}[целочисленная арифметика и платёжеспособность при округлении]
Теорема~\ref{thm:lp} сформулирована над рациональными числами, но протокол
оперирует целыми числами (милли-VIZ), и каждое деление применяет оператор пола
$\lfloor \cdot \rfloor$. Покажем, что гарантия переживает дискретизацию. Каждая
комиссия и каждая доля выигравшего округляются \emph{вниз}:

\[
f_\bullet = \lfloor S_{\text{lose}} \cdot \theta_\bullet / 10000 \rfloor, \qquad
\hat{\pi}_i^{\text{profit}} = \left\lfloor \frac{w_i}{\sum_{j} w_j} \cdot W
\right\rfloor
\]

Поскольку $\lfloor x \rfloor \leq x$, фактическое распределение выигравшим
удовлетворяет $\sum_{i \in \mathcal{W}} \hat{\pi}_i^{\text{profit}} \leq W$, и
аналогично $\sum_\bullet f_\bullet \leq S_{\text{lose}} \cdot (\theta_{\text{oracle}} +
\theta_{\text{creator}} + \theta_{\text{liq}})/10000 \leq S_{\text{lose}}$. Поэтому
фактический отток подчиняется

\[
\text{Money}_{\text{OUT}}^{\text{realized}} = L + \sum_{i \in \mathcal{W}} a_i +
\sum_{i \in \mathcal{W}} \hat{\pi}_i^{\text{profit}} + \sum_\bullet f_\bullet
\;\leq\; L + \sum_{\text{all bets}} a_i = \text{Money}_{\text{IN}}
\]

так что система никогда не становится неплатёжеспособной. Неотрицательный остаток

\[
\epsilon = S_{\text{lose}} - \sum_{i \in \mathcal{W}}
\hat{\pi}_i^{\text{profit}} - \sum_\bullet f_\bullet \;\geq\; 0
\]

ограничен $\epsilon < |\mathcal{W}| + 3$ милли-VIZ (одна субъединица на каждую
округлённую величину) и сметается в фонд DAO как «пыль». Таким образом, округление
может лишь \emph{недо}-распределять, но никогда не пере-распределять: основной
капитал LP $L$ сохраняется точно, а неравенство $\text{Money}_{\text{OUT}} \leq
\text{Money}_{\text{IN}}$ выполняется на гранулярности реестра.
\end{remark}

\subsection{Временной штраф за поздние ставки}\label{sec:timepenalty}

Чтобы смягчить информационную асимметрию ставок вблизи истечения (когда
неопределённость исхода снижена), протокол вводит настраиваемый временной штраф,
применяемый \emph{только к прибыли}, никогда --- к основному капиталу.

Пусть $\Delta t = t_{\text{bet}} - t_{\text{now}}$ --- время до истечения, а
$\omega$ --- окно штрафа. Если $\Delta t < \omega$:

\[
r = 1 - \frac{\Delta t}{\omega}
\]

\[
\tau_{\text{rate}} = \begin{cases} r & \text{линейный} \\ r^2 &
\text{квадратичный (по умолчанию)} \end{cases}
\]

\[
\tau_i = \lfloor \pi_{\text{profit},i} \cdot \tau_{\text{rate}} \cdot
\tau_{\max} / 10^6 \rfloor
\]

где $\pi_{\text{profit},i}$ --- доля прибыли беттора (без основного капитала), а
$\tau_{\max}$ --- максимальный штраф в микроединицах.

\textbf{Инвариант:} $\pi_i \geq a_i$ для всех выигравших бетторов, поскольку штраф
применяется исключительно к компоненте прибыли.

\subsection{Взвешенное по времени распределение комиссий LP}\label{sec:lpfee}

Доли комиссий LP распределяются пропорционально произведению суммы вклада на
остаточное время до истечения:

\[
\omega_j = \text{amount}_j \cdot \max(1, t_{\text{bet}} - t_{\text{deposit},j})
\]

\[
f_{\text{share},j} = \left\lfloor \frac{f_{\text{pool}} \cdot \omega_j}{\sum_m
\omega_m} \right\rfloor
\]

Это создаёт сильный стимул к раннему предоставлению ликвидности. На рынке
длительностью $T$ LP, вносящий депозит в момент $t = 0$, зарабатывает в $T$ раз
больше на единицу капитала, чем вносящий в $t = T - 1$.

% =====================================================================
\section{«Ленивый» пул ликвидности}\label{sec:lazypool}

\subsection{Мотивация}

Индивидуальное предоставление ликвидности требует активной оценки и выбора рынков.
Для розничных участников это создаёт непрактичный барьер знаний и усилий.
«Ленивый» пул (так названный потому, что не требует от вкладчиков активного выбора
рынков) решает это, принимая депозиты и автоматически размещая капитал в
активируемые рынки.

\subsection{Механика пула}\label{sec:poolmech}

«Ленивый» пул $\mathcal{L}_P$ --- это синглтон-контракт, поддерживающий:
\begin{itemize}
  \item $B_{\text{free}}$: нераспределённый баланс
  \item $B_{\text{alloc}}$: капитал, размещённый в активных рынках
  \item $B_{\text{earned}}$: реализованная прибыль
  \item $S_{\text{total}}$: общее число выпущенных долей
  \item $\rho$: накопитель совокупного вознаграждения на долю
\end{itemize}

\textbf{Депозит.} Когда пользователь $u$ вносит сумму $a$:

\[
\text{shares}_u = \begin{cases} a & \text{если } S_{\text{total}} = 0 \\
a \cdot S_{\text{total}} / B_{\text{free}} & \text{иначе} \end{cases}
\]

Депозит блокируется на заданный управлением период $t_{\text{lock}}$.

\textbf{Авто-размещение.} При активации рынка ($q_0 \rightarrow q_1$):

\[
a_{\text{alloc}} = B_{\text{free}} \cdot \alpha / 100
\]

при условиях $a_{\text{alloc}} \geq a_{\min}$ и $B_{\text{alloc}} +
a_{\text{alloc}} \leq B_{\text{total}} \cdot \alpha_{\max} / 100$, где $\alpha$ ---
процент размещения на рынок, а $\alpha_{\max}$ --- потолок суммарного размещения.

Размещение дополнительно ограничено \textbf{порогом вознаграждения}: пул
предоставляет капитал рынку, только если его LP-комиссия $f_{\text{liq}}$ достигает
управляемого минимума $f_{\text{liq}}^{\min}$ (по умолчанию $2\%$). Рынок,
предлагающий LP меньше этого, не окупает альтернативные издержки пула, поэтому пул
размещает ноль, и единственная глубина рынка --- собственный сид создателя. Это
делает LP-комиссию явной ценой, которую рынок платит за автоматическую пуловую
ликвидность, и не даёт пулу субсидировать рынки, настроенные направлять всю
комиссионную выручку мимо поставщиков ликвидности.

Формула размещения включает поправки на качество оракула:

\[
a_{\text{alloc}} = B_{\text{free}} \cdot \frac{\alpha}{100} \cdot
(1 - \beta)^{n_o} \cdot (1 - \gamma)^{f_o}
\]

где $n_o$ --- число активных рынков оракула, $f_o$ --- число активных штрафных
меток оракула, а $\beta, \gamma$ --- заданные управлением коэффициенты затухания.

\subsection{Распределение вознаграждений («ленивый» учёт)}\label{sec:lazyacct}

Вознаграждения распределяются с помощью единственного глобального накопителя
$\rho$ (вознаграждение на долю) по паттерну MasterChef (SushiSwap, 2020; Leshner \&
Hayes, 2019):

Когда рынок разрешается и LP-позиция пула приносит прибыль $\pi_{\text{pool}}$:

\[
\rho \leftarrow \rho + \frac{\pi_{\text{pool}} \cdot \text{PRECISION}}{S_{\text{total}}}
\]

Накопленное вознаграждение любого пользователя $u$ вычислимо за $O(1)$:

\[
\text{reward}_u = \text{pending}_u + \frac{\text{shares}_u \cdot (\rho -
\rho_{\text{snapshot},u})}{\text{PRECISION}}
\]

Записи пользователя обновляются только при его действии (депозит или вывод), что
даёт амортизированную стоимость $O(1)$ независимо от числа участников.

\subsection{Защита от издержек упущенных возможностей}\label{sec:oppcost}

Механизм авто-размещения создаёт вектор атаки: злонамеренный оракул может создавать
долгие рынки с нулевым объёмом, чтобы заблокировать капитал пула. Этому
противодействуют три градуированных защитных механизма:

\textbf{Градуированный отзыв.} Длительность рынка делится на $n_{\text{steps}}$
интервалов. На каждой контрольной точке $k$, если накопленный объём $V_k$
удовлетворяет:

\[
V_k < a_{\text{alloc}} \cdot \theta_{\text{recall}} / 100
\]

то доля $\delta_{\text{recall}}$ текущего размещения отзывается в пул. После $n$
подряд идущих низкообъёмных контрольных точек удержанное размещение равно:

\[
a_{\text{retained}} = a_{\text{alloc}} \cdot (1 - \delta_{\text{recall}})^n
\]

Для параметров по умолчанию ($\delta_{\text{recall}} = 10\%$, $n_{\text{steps}} =
10$) полностью простаивающий рынок удерживает $(0.9)^{10} \approx 35\%$ исходного
размещения.

\textbf{Штраф за активные рынки.} Каждый активный рынок оракула $o$
мультипликативно снижает новые размещения на множитель $(1 - \beta)$. При $\beta =
5\%$ и 10 активных рынках новые размещения составляют $(0.95)^{10} \approx 60\%$
базовой ставки.

\textbf{Штрафные метки.} Плохие исходы (no-contest, пропуск дедлайнов, проигранные
споры, разрешения с нулевым объёмом) порождают штрафные метки на оракуле,
дополнительно снижающие размещения на множитель $(1 - \gamma)^{f_o}$. Метки
автоматически истекают после заданного управлением окна чистой работы.

\subsection{Экстренный вывод}

Пользователи могут выйти из пула до истечения блокировки, при штрафе \emph{только
на заблокированную прибыль}:

\[
\text{penalty} = \max(0, V_{\text{total}} - P_{\text{deposited}}) \cdot
\frac{S_{\text{locked}}}{S_{\text{total}}} \cdot \theta_{\text{emergency}} / 100
\]

где $V_{\text{total}}$ --- общая стоимость пользователя (доли + накопленные
вознаграждения), $P_{\text{deposited}}$ --- совокупно внесённый основной капитал, а
$S_{\text{locked}} / S_{\text{total}}$ --- доля ещё заблокированных долей. Штраф
перераспределяется оставшимся участникам пула через $\rho$.

\subsection{Опциональное кредитное плечо, финансируемое из пула}\label{sec:leverage}

«Ленивый» пул играет вторую, опциональную роль из того же $B_{\text{free}}$: он
финансирует подсистему кредитного плеча. Эта подсистема \textbf{опциональна,
выключена по умолчанию и управляется медианным аварийным выключателем
(kill-switch)}. Мы рассматриваем её отдельно от остальной статьи по важной причине.
\textbf{Это единственное место, где капитал пула принимает на себя кредитный риск,
поэтому оно \emph{не} покрывается безусловной гарантией Теоремы~\ref{thm:lp}.}
Теорема~\ref{thm:lp} касается пула в роли \emph{поставщика рыночной ликвидности},
где основной капитал структурно изолирован от исходов ставок. Кредитное плечо же
ставит пул в роль \emph{кредитора}, а кредитование несёт риск, который мы
ограничиваем, а не устраняем.

\textbf{Открытие позиции.} Беттор вносит залог $m$. Пул ссужает маржу $\lambda m$
при коэффициенте плеча $\lambda$ из $B_{\text{free}}$, при потолках на размер
позиции, долю фонда пула и минимальную ликвидность рынка. Ни один токен не
создаётся: суммарная ставка $(1+\lambda)m$ входит в рынок точно так же, как обычная
ставка. Вес кривой позиции, её залог и ссуженная сумма записываются в выделенный
объект позиции. Обязательство пула по позиции:

\[
\Omega = \lambda m \,(1 + r),
\]

где $r$ --- накопленный процент. Ссуда плюс процент --- это то, что пул стремится
вернуть.

\textbf{Ликвидация против записанного снимка резервов.} Трудный случай на бинарных
рынках с плечом --- \emph{скачковый риск} (jump risk): разрывное движение цены может
оставить наивную ликвидацию «продать по текущей цене» неспособной покрыть ссуду.
Onix не ликвидирует по текущей цене. Он закрывает позицию через обратный CPMM
против \textbf{состояния резервов, записанного для этой позиции}, так что
собственное рыночное воздействие позиции разматывается первым, а не оплачивается
дважды. Возврат обеспечивают два пути:

\begin{enumerate}
  \item \textbf{Каскад встречной ставки} (\texttt{pm\_place\_bet}, причина
  ликвидации 0) --- входящая встречная ставка запускает принудительное закрытие
  позиции с плечом против её записанного снимка.
  \item \textbf{Принудительное закрытие при расчёте} (\texttt{pm\_leverage\_resolve})
  --- при разрешении позиция закрывается, и ссуда погашается из её веса прежде, чем
  заёмщику начислится прибыль.
\end{enumerate}

На этих путях проектная цель --- полный возврат ссуды, $\text{recovered} =
\min(V_{\text{close}}, \Omega) \ge \lambda m$, с возвратом фактической ссуды и
процента в пул. Мы формулируем это как \textbf{проектное свойство, а не теорему}:
математика расчёта по плечу реализована и проверяется в \texttt{consensus\_sim}, но
замкнутое доказательство возврата при всех траекториях резервов оставлено на
будущую работу (\Cref{sec:limitations}).

\textbf{Куда теперь идёт апсайд заёмщика.} Поскольку плечо закрывается против кривой,
а не как тотализаторная доля, его выход управляется механизмом отложенного
требования при раннем выходе (\Cref{sec:earlyexit}): пул сначала забирает своё
обязательство $\Omega$ из значения закрытия, а любой остаток $\max(V_{\text{close}}
- \Omega, 0)$ становится \emph{отложенным правом требования, зависящим от исхода}, а
не немедленной выплатой. Плечо, таким образом, есть \textbf{направленная ставка с
плечом}, а не сбор волатильности --- заёмщик получает прибыль, только если выбранный
исход побеждает и в капнутом бакете есть место. Именно это закрывает исторический
путь непокрытой недостачи (\Cref{sec:earlyexit}); с ним апсайд заёмщика никогда не
может черпаться из основного капитала LP.

\textbf{Аварийный выключатель отвязан от защиты.} Аварийный выключатель отключает
только \emph{новые} открытия. Пути ликвидации намеренно \textbf{не} ограничены им,
поэтому отключение плеча никогда не лишает защиты уже открытые позиции.

\textbf{Изоляция риска и вознаграждение.} Кредитование с плечом ограничено
управляемой долей $B_{\text{free}}$ (\texttt{leverage\_fund\_percent}), так что
экспозиция пула в худшем случае ограничена на уровне подсистемы и не может достичь
основного капитала рыночного LP, защищённого Теоремой~\ref{thm:lp}. На пути
добровольной отмены той же стороны заёмщик, закрывающийся при неблагоприятном
движении, всё ещё может оставить пул с недостачей до $\Omega$; эта недостача
\textbf{ограничена залогом заёмщика $m$}, который протокол удерживает как маржу
первого убытка, и не может его превысить. Заработанный процент начисляется в пул
через тот же накопитель $\rho$ (\Cref{sec:lazyacct}). Поэтому вкладчики пула
зарабатывают из двух источников --- доли комиссий из пула проигравших (без риска,
Теорема~\ref{thm:lp}) и процент по плечу (ограниченный кредитный риск, опционально)
--- и оба учитываются одинаково, но управляются независимо.

\subsection{Дизайн-решение: только реальная глубина (без виртуальной/фантомной ликвидности)}\label{sec:realdepth}

Естественное предложение --- засеять кривую ценообразования рынка \emph{виртуальной}
(фантомной) ликвидностью: оффсетом резерва $\phi$, который добавляют, чтобы уплостить
влияние на цену, но который не обеспечен реальным капиталом и удаляется при
расчёте, опционально настраиваемым управлением. В закрытом цикле «ставка
$\rightarrow$ отмена $\rightarrow$ расчёт» это сохраняет стоимость (это техника
виртуального AMM) и соблазнительно как «стабилизатор» холодного старта для новых
тонких рынков. Onix \textbf{сознательно этого не внедряет.} Ту же пользу на старте
уже даёт авто-размещение «ленивого» пула (\Cref{sec:poolmech}) --- но \emph{реальным}
капиталом, который вдобавок зарабатывает комиссии, имеет подотчётного владельца и
следует за спросом по каждому рынку. Мы отвергаем фантомную глубину, потому что при
неосторожном применении она вредит ровно тому, что протокол призван защищать, ---
структуре рынка и доверию:

\begin{enumerate}
  \item \textbf{Подделываемая глубина подрывает доверие.} Глубина значима как
  сигнал лишь потому, что это дорогой реальный капитал под риском. Бесплатная
  виртуальная глубина превращает заявление «рынок глубокий и ликвидный» в
  подделываемое --- тонкий или манипулируемый рынок можно нарядить под глубокий.
  Реальный капитал делает это заявление неподделываемым.
  \item \textbf{Она тихо искажает агрегирование информации.} Виртуальная глубина
  уплощает кривую весов, ослабляя награду за раннюю верную информацию и делая
  отображаемую цену невосприимчивой к новостям («стабильна, потому что её нельзя
  сдвинуть» $=$ устаревший прогноз). Правильная величина зависит от рынка и объёма;
  единственная константа управления не может за ней следовать и, заданная слишком
  высоко, разрушает то самое свойство открытия цены, ради которого существует рынок
  предсказаний.
  \item \textbf{Она платёжеспособна, только пока её не выкупают и не используют как
  залог.} Как только глубина обеспечивает реальный отток --- отмены, ранние выводы,
  ссуды плеча, общие/межрыночные пулы, --- виртуальную часть приходится исключать
  везде, иначе она утекает реальными деньгами (например, ссуда плеча, рассчитанная
  или возвращаемая против фейковой глубины, превращается в реальный безнадёжный долг
  вкладчикам пула). Каждая ветка «считать против ликвидности / платить LP»
  становится граблями, требующими «\ldots но не фантомную часть». Реальный капитал
  устраняет весь этот класс ошибок исключаемости по построению.
  \item \textbf{У неё нет владельца, дохода и подотчётности.} Виртуальная глубина
  не несёт риска и не зарабатывает комиссию никому реальному; это услуга, за которую
  никто не платит и за которую никто не отвечает. На рынках, которых она касается,
  она удаляет розничный продукт безрисковой доходности --- основную гипотезу
  протокола (Q3, H1).
\end{enumerate}

«Ленивый» пул --- та же идея, сделанная на реальных числах: авто-размещение
сглаживает запуск новых рынков, но капитал погашаем, безопасен для плеча,
зарабатывает комиссии, принадлежит вкладчикам и самокорректируется по каждому рынку
через градуированный отзыв (\Cref{sec:oppcost}). Поэтому Onix держит \textbf{только
реальные числа} --- каждая единица глубины есть реальный капитал, который можно
вывести, который зарабатывает и который подотчётен. Это осознанный компромисс: мы
отказываемся от дешёвого виртуального стабилизатора ради целостности ценового
сигнала и платёжеспособности каждого реального денежного пути.

% =====================================================================
\section{Ранний выход через отложенные права требования, зависящие от исхода}\label{sec:earlyexit}

\subsection{Напряжение «выход по кривой / тотализатор»}\label{sec:exittension}

Рынок Onix гибриден: \textbf{кривая} CPMM (бинарный) / LMSR (многоисходный)
ценообразует и вход, и ранний выход, тогда как позиции, \emph{удерживаемые до
разрешения}, рассчитываются тотализаторно (\Cref{sec:settlement}). Это создаёт
структурное напряжение. Любой круговой проход через кривую --- купить, затем продать
до расчёта --- реализует торговую P\&L против глубины кривой, то есть против LP,
ровно как непостоянные потери Uniswap. Но Теорема~\ref{thm:lp} обещает LP
\emph{защиту основного капитала} с апсайдом только от комиссий. Эти два свойства
прямо конфликтуют на любом пути раннего выхода.

Два консенсусных пути выходят против кривой на бинарном рынке:

\begin{itemize}
  \item \textbf{Плечо} (\texttt{liquidate\_position} / \texttt{pm\_leverage\_close})
  --- всегда, и с принудительным закрытием при расчёте, усиленное ссудой пула
  (\Cref{sec:leverage}).
  \item \textbf{Отмена обычной ставки} (\texttt{cancel\_bet}, возврат по цене кривой)
  --- в течение окна приёма ставок.
\end{itemize}

В наивном дизайне оба пути направляют знаковый остаток
$\text{residual} = \text{stake} - \text{curve\_refund}$ в $\text{forfeit\_pool}$
рынка. Когда ранний выход \emph{прибылен} ($\text{curve\_refund} > \text{stake}$),
остаток отрицателен и $\text{forfeit\_pool}$ становится отрицательным. При расчёте
пул победителей равен

\[
\text{winners\_pool} = S_{\text{lose}} - \text{fees} + \text{forfeit\_pool},
\]

так что если прибыль раннего выхода (особенно с плечом) превысила ставки
проигравших, $\text{winners\_pool} < 0$. Прижатая полом к нулю, недостача --- назовём
её $\text{uncovered}$ --- списывается на основной капитал LP, а как только основной
капитал LP исчерпан, фактически \emph{эмитируется}. Это не гипотетично: она
достижима в соблюдающей гейты CPMM-симуляции (односторонний «памп» даёт
$\text{uncovered} = 7316$ милли-VIZ), и состязательный replay-корпус попал в неё в
$1163/1988$ парных траекторий.

Корневая причина точна: \textbf{выход по цене кривой платит значение бондинг-кривой,
которое не ограничено пулом проигравших}, тогда как расчёт платит тотализаторно
(ограничен $S_{\text{lose}}$). Неограниченный разрыв ложится на LP --- ломая ту самую
гарантию, ради которой существует архитектура.

\subsection{Механизм отложенного требования}\label{sec:defclaim}

Onix разрешает это напряжение, делая так, чтобы ранние выходы \emph{вообще не}
извлекали значение кривой у LP. Вместо этого выход записывает \textbf{отложенное
право требования, зависящее от исхода}, финансируемое при расчёте из
\emph{ограниченного слайса пула проигравших}.

\textbf{Записывается при выходе.} Объект отложенного требования хранит
$\big(\text{position\_id},\ \text{kind} \in \{\text{bet}, \text{leverage}\},\
\text{chosen\_outcome},\ \text{claim\_amount},\ \text{exit\_time}\big)$.

\textbf{Отмена обычной ставки.} Основной капитал возвращается \emph{немедленно и
безусловно},
\[
\text{refund} = \min(\text{curve\_refund}, \text{stake}),
\]
потому что это собственные деньги беттора. Отмена может срезать убытки или выйти в
ноль, но никогда не реализует прибыль кривой в момент отмены. Хвост прибыли
\[
\text{claim\_amount} = \max(\text{curve\_refund} - \text{stake}, 0)
\]
становится отложенным требованием на выбранный исход. Критически, только
\emph{неотрицательный} остаток $\text{stake} - \text{refund} \ge 0$ вообще
направляется в $\text{forfeit\_pool}$; путь отрицательного остатка, порождавший
$\text{uncovered}$, устранён.

\textbf{Закрытие / ликвидация плеча.} Залог \emph{не} возвращается отдельно --- это
маржа первого убытка для пула. Пул сначала возвращает своё обязательство
$\Omega = \lambda m (1 + r)$ из значения закрытия $V_{\text{close}}$; если
$V_{\text{close}} < \Omega$, залог покрывает разрыв. Остаток
\[
\text{claim\_amount} = \max(V_{\text{close}} - \Omega, 0)
\]
становится отложенным требованием на выбранный исход. Позицией с плечом теперь
управляют два разных предиката, и их нельзя смешивать: \emph{платёжеспособность}
($V_{\text{close}} \ge \Omega$, определяющая возврат ссуды пулу) против
\emph{победы исхода} (определяющей право на требование). Позиция может быть
платёжеспособной, но на проигравшем исходе --- пул получает своё сполна, а
требование равно нулю.

\subsection{Распределение при расчёте}\label{sec:defsettle}

При расчёте отложенные требования финансируются из капнутого бакета:

\begin{enumerate}
  \item Вычислить бакет, ограниченный (клэмп) свободным запасом расчёта $H$
  (headroom, \Cref{sec:defthm}), так что он никогда не может превысить пул,
  доступный победителям:
  \[
  \text{bucket} = \min\!\left(\left\lfloor
  \frac{\texttt{pm\_early\_exit\_reward\_cap\_percent} \cdot S_{\text{lose}}}{10000}
  \right\rfloor,\; H\right)
  \]
  (кап по умолчанию $=3300$ bp $=33\%$ пула проигравших; ограничение связывает лишь
  тогда, когда комиссии достаточно велики, чтобы кап иначе перечерпал бы пул).
  \item Собрать отложенные требования \textbf{только на выигравшем исходе};
  требования на проигравшем исходе платят ноль.
  \item Платить их \textbf{FIFO по \text{exit\_time}} (кто раньше вышел, тому раньше
  платят), пока бакет не опустеет. Пер-позиционного капа нет --- FIFO-порядок плюс
  суммарный бакет \emph{и есть} граница. Требование, которое остаток бакета не может
  профинансировать полностью, оплачивается частично; непрофинансированный остаток ---
  это урезание (haircut).
  \item \textbf{Любой неиспользованный бакет возвращается в пул победителей}, где
  удерживаемые выигравшие ставки делят его тотализаторно.
\end{enumerate}

\subsection{Теорема: ограниченный ранний выход, без непокрытой недостачи}\label{sec:defthm}

\begin{theorem}[Ограниченный ранний выход]\label{thm:earlyexit}
Пусть $c = \texttt{pm\_early\_exit\_reward\_cap\_percent}/10000 \in [0,1)$, и пусть
финансирующий бакет ограничен свободным запасом расчёта:
\[
\text{bucket} = \min\!\big(\lfloor c\, S_{\text{lose}}\rfloor,\; H\big), \qquad
H = \max\!\big(0,\; S_{\text{lose}} - \text{fees} - f_{\text{fixed}}^{\text{paid}}
+ \text{honest\_forfeits}\big),
\]
где $H$ --- пул победителей, доступный непосредственно \emph{перед} оплатой
требований, и $\text{honest\_forfeits} \ge 0$. Тогда для \emph{любого} рынка и
\emph{любых} валидных параметров комиссий оплаченные требования раннего выхода
никогда не превышают бакет, пул победителей неотрицателен, и основной капитал LP
никогда не затрагивается ранними выходами:
\[
\text{paid\_claims} \;\le\; \text{bucket} \;\le\; H, \qquad
\text{winners\_pool} \;=\; H - \text{paid\_claims} \;\ge\; 0.
\]
\end{theorem}

\begin{proof}
FIFO-распределение (\Cref{sec:defsettle}) прекращает платить, как только
накопленные выплаты достигают бакета, поэтому $\text{paid\_claims} \le
\text{bucket}$. По ограничению $\text{bucket} \le H$, следовательно
$\text{paid\_claims} \le H$. Пул победителей равен ровно $H$ минус то, что выплачено
ранним выходящим, поэтому $\text{winners\_pool} = H - \text{paid\_claims} \ge 0$.
Никакая недостача не списывается на основной капитал LP, и поскольку отток никогда
не превышает $S_{\text{lose}} + L + \sum_{\mathcal{W}} a_i$, токены не эмитируются.
После механизма \Cref{sec:defclaim} только неотрицательные остатки достигают
$\text{forfeit\_pool}$, поэтому $\text{honest\_forfeits} \ge 0$ и $H$ корректно
определён.
\end{proof}

Именно ограничение (clamp) делает гарантию \emph{безусловной}. Без него
понадобился бы управленческий инвариант $\theta_{\text{total}}/10000 + c \le 1$ (с
$\theta_{\text{total}} = \theta_{\text{oracle}} + \theta_{\text{creator}} +
\theta_{\text{liq}}$), дающий интуитивную границу $\text{winners\_pool} \ge
(1-c)\,S_{\text{lose}} - \text{fees} \ge 0$. Этот инвариант \emph{не} обеспечивается
валидацией параметров --- пер-рыночные комиссии ограничены лишь условием
$\theta_{\text{total}} \le 100\%$ при создании рынка, без потолка на уровне цепи на
комиссию создателя или комиссию ликвидности, --- так что валидный рынок с высокими
комиссиями в сочетании с ненулевым капом мог бы иначе увести $\text{winners\_pool}$
в отрицательную область и списать недостачу на основной капитал LP. В режиме по
умолчанию (комиссии $\approx 2\%$, кап $33\%$) $\min$ неактивен и $\text{bucket} =
\lfloor c\,S_{\text{lose}}\rfloor$, точно восстанавливая задуманный кап; ограничение
связывает лишь в состязательном углу с высокими комиссиями, где оно сохраняет
платёжеспособность ценой более глубокого урезания раннего выхода.

\begin{corollary}[без эмиссии, безопасно для LP, проигравший никогда не в прибыли]
При условиях Теоремы~\ref{thm:earlyexit}: (i) недостача $\text{uncovered}$
невозможна по построению, поэтому эмиссии не происходит, а основной капитал LP ---
включая основной капитал рыночного LP «ленивого» пула --- никогда не затрагивается
ранними выходами или непостоянными потерями плеча; (ii) требование на
\emph{проигравшем} исходе платит ноль, поэтому проигравший исход никогда не бывает в
прибыли; и (iii) удерживаемые победители получают не менее $(1-c)\,S_{\text{lose}}$
пула проигравших плюс любой неиспользованный бакет.
\end{corollary}

\subsection{Экономическая интерпретация: ограниченный дисконт за ликвидность}\label{sec:defecon}

Отложенное требование переформулирует ранний выход как \emph{выбор} с чистыми
стимулами. Держатель, остающийся до разрешения, получает полную тотализаторную долю
пула проигравших. Держатель, выходящий рано, получает вместо этого (a) немедленный
возврат основного капитала (обычная ставка) или обработку залога с погашенной ссудой
(плечо), плюс (b) \emph{капнутое, приоритетно упорядоченное, зависящее от исхода}
требование не более чем на $c$ пула проигравших, оплачиваемое FIFO. Поскольку выплата
раннего выхода строго является капнутой-и-зависящей-от-исхода версией выплаты за
удержание, \textbf{арбитража против удержания нет}: нельзя выйти рано, чтобы извлечь
\emph{больше}, чем принесло бы удержание. Дисконт намеренный --- это цена ликвидности
и цена более раннего переноса риска исхода, --- и именно он позволяет гарантии LP
выжить при кривой, которая также ценообразует выходы.

Это также придаёт плечу связную экономическую идентичность. По правилу отложенного
требования заём из «ленивого» пула для открытия позиции --- это \emph{направленная}
ставка с плечом: заёмщик усиливает экспозицию на выбранный исход, пулу сначала
возвращаются ссуда и процент из значения закрытия по кривой, а апсайд заёмщика ---
ограниченное зависящее от исхода требование, которое конкурирует FIFO с другими
ранними выходящими за капнутый бакет. Интерпретация «сбор волатильности» --- открыть,
прокатиться по кривой, обналичить колебание за счёт LP --- структурно устранена.

\subsection{Параметр управления и реализация}\label{sec:defimpl}

Кап --- это выбираемый медианным голосованием валидаторов параметр
\texttt{pm\_early\_exit\_reward\_cap\_percent} (uint16, базисные пункты, по умолчанию
$3300$), добавленный в \texttt{chain\_properties\_pm}, причём $\texttt{validate()}$
ограничивает его значением $\le 10000$. Отложенные требования живут в выделенном
объекте \texttt{pm\_deferred\_claim\_object} (индексированном по рынку и по порядку
выхода), включённом в allowlist снапшота с импорт-обработчиком. В цикле распределения
при расчёте эмитируется виртуальная операция
\texttt{pm\_early\_exit\_claim\_paid}$(\text{account}, \text{market},
\text{kind}, \text{outcome}, \text{claimed}, \text{paid})$ --- раскрывающая любое
урезание из-за исчерпания бакета через разрыв между \texttt{claimed} и \texttt{paid}
--- и направляется в историю аккаунта выходящего (голая корректировка баланса не
оставила бы следа). Read-API \texttt{get\_deferred\_claims} возвращает FIFO-очередь,
пустую на рассчитанном рынке после потребления требований.

Три защитные меры укрепляют механизм против состязательных экстремумов параметров и
отказа в обслуживании. Во-первых и важнее всего, финансирующий бакет
\textbf{ограничен свободным запасом расчёта} $H$ (headroom, \Cref{sec:defthm}):
сумма, фактически черпаемая на требования, равна $\min(\lfloor c\,S_{\text{lose}}\rfloor, H)$,
что и делает Теорему~\ref{thm:earlyexit} безусловной, а не зависящей от
управленческого инварианта $\text{fees} + \text{cap} \le 100\%$ --- последний не
обеспечивается валидацией параметров, поскольку у комиссий создателя и ликвидности
нет потолка на уровне цепи. Во-вторых, пер-рыночное число объектов отложенных
требований ограничено \texttt{MAX\_PM\_DEFERRED\_CLAIMS\_PER\_MARKET}; сверх капа
ранний выход пропускает создание требования, и хвост просто остаётся в кривой, платя
ноль --- он намеренно \emph{не} направляется в $\text{forfeit\_pool}$, что привело бы
к двойному учёту и, следовательно, к эмиссии. В-третьих, всегда включённая
диагностика срабатывает, если когда-либо возникает $\text{uncovered} > 0$
(недостижимо после ограничения бакета), а пол $\lfloor \cdot \rfloor_{\ge 0}$ на пуле
победителей остаётся как последний рубеж. Это двойная подстраховка: ограничение даёт
гарантию по построению, а диагностика делает любое латентное нарушение громким, а не
тихим.

\textbf{Верификация.} Инвариант сохранения из Теоремы~\ref{thm:earlyexit} был
проверен в модели расчёта по состязательным сценариям (односторонние «пампы», тонкие
проигравшие стороны, исчерпание бакета, экстремумы комиссий и
$\Omega > \text{суммарная ставка}$ от долгоживущего плеча) и в харнессе
\texttt{consensus\_sim}. Он был далее подтверждён сквозным прогоном на тестнете VIZ:
отмена записала отложенное требование; при расчёте виртуальная операция сообщила о
выплате требования из капнутого бакета с видимым урезанием
($\texttt{claimed} > \texttt{paid}$), и требование было потреблено --- полный
жизненный цикл вёл себя ровно так, как предсказывает модель.

% =====================================================================
\section{Разрешение споров под управлением DAO}\label{sec:dispute}

\subsection{Модель оракула с залогом}

Каждый оракул $o$ обязан поддерживать страховой залог $I_o \geq I_{\min}$. Залог
создаёт подотчётность: у оракулов, которые неверно разрешают рынки, пропускают
дедлайны или проигрывают споры, страховка списывается. Доход оракула складывается
из:

\begin{enumerate}
  \item Фиксированной комиссии за рынок $f_{\text{fixed}}$ (компенсирующей внесение
  страховки)
  \item Процентной комиссии $f_{\text{oracle}}$ из пула проигравших при разрешении
\end{enumerate}

Эти две комиссии входят в систему в разных точках, и их нельзя смешивать.
\textbf{Процентная комиссия} $f_{\text{oracle}}$ удерживается из $S_{\text{lose}}$
при разрешении и фигурирует в денежном потоке Теоремы~\ref{thm:lp}
(\Cref{sec:thm1}). \textbf{Фиксированная комиссия} $f_{\text{fixed}}$ --- это прямой
перевод создатель$\rightarrow$оракул, выполняемый при \emph{принятии} рынка, до
каких-либо ставок; поэтому она \emph{не} является частью денежного потока
ставок/LP --- она не добавляется к начальной ликвидности $L$ и не черпается из
$S_{\text{lose}}$, а значит не появляется в Теореме~\ref{thm:lp}. (Она опущена в
векторе комиссий $\boldsymbol{\theta}$ из \Cref{sec:fees} по той же причине:
$\boldsymbol{\theta}$ собирает только процентные комиссии на момент разрешения,
финансируемые проигравшими.) Для рынков с самооракулом перевод не происходит и
$f_{\text{fixed}} = 0$.

Процесс принятия оракулом реализует механизм «оферта--котировка»: создатель
публикует потолки комиссий, а оракул фиксирует свои фактические условия
(ограниченные как потолком создателя, так и потолком управления) при принятии.
Принятие ограничено во времени окном $\tau_{\text{accept}}$ (\Cref{sec:fsm}): рынок,
оставшийся в ожидании дольше этого срока, авто-аннулируется поблочным кроном, который
возвращает начальную ликвидность создателя, удерживая антиспам-комиссию за создание,
--- так неотвечающий оракул не может бесконечно замораживать сид.

\subsection{Двухрежимная система споров}\label{sec:disputemodes}

Любой беттор может оспорить разрешение в течение льготного периода
$\Delta t_{\text{grace}}$, внеся в эскроу комиссию спора $d$. Во время споров все
выплаты заморожены. Протокол поддерживает два режима спора на рынок:

\textbf{Режим комитета ($\delta_{\text{mode}} = 0$).} Весь электорат держателей
токена разрешает спор взвешенным по стейку публичным голосованием. Вес каждого
голосующего $v$:

\[
w_v = s_v + \text{shares}_{v,\text{pool}} \cdot \text{NAV}_{\text{pool}} /
S_{\text{total}}
\]

где $s_v$ --- вестинг-доли голосующего, а второе слагаемое переводит стейк в
«ленивом» пуле в эквивалентный вес управления. Голоса публичны и пересматриваемы до
закрытия периода голосования --- это намеренное проектное решение: споры суть
\emph{прозрачные публичные слушания}, а не тайные бюллетени. Новые данные могут
изменить голоса.

Эта прозрачность несёт известную цену. Поскольку текущий подсчёт виден, голосующий
может отложить голос на конец окна и обусловить свой бюллетень раскрытыми позициями
других (преимущество последнего хода). Протокол принимает этот компромисс намеренно:
для DAO аудируемость и легитимность открытого слушания ценнее тайны бюллетеня, а
переворот исхода всё равно требует сдвинуть взвешенное по стейку большинство.
Манипуляция сдерживается не тайной, а (i) взвешенным по стейку порогом одобрения
$\theta_{\text{approve}}$, (ii) многодневным окном голосования, размывающим любое
единичное преимущество тайминга, и (iii) согласованием вознаграждения арбитра,
анализируемым в \Cref{sec:gametheory}. Бюллетень commit-reveal был явно рассмотрен и
\textbf{отвергнут} по этой причине --- сокрытие голосов разрушило бы свойство
публичного слушания, дающее вердикту легитимность. Мы возвращаемся к этому как к
честному ограничению в \Cref{sec:limitations}.

Спор удовлетворяется, если процент одобрения превышает заданный управлением порог
$\theta_{\text{approve}}$:

\[
\frac{\sum_{v: \text{vote}_v = \text{approve}} w_v}{\sum_{v} w_v} \geq
\frac{\theta_{\text{approve}}}{10000}
\]

\textbf{Режим аккаунта ($\delta_{\text{mode}} = 1$).} Именованный арбитр спора
(рекомендуется: мультиподпись) выносит обязывающий вердикт.

\subsection{Исходы спора и совместимость стимулов}

\textbf{Оракул неправ (спор удовлетворён):}

\[
\text{reward}_{\text{pool}} = \min(d \cdot \mu, I_o)
\]

где $\mu$ --- множитель вознаграждения. Пул вознаграждения делится так:

\begin{itemize}
  \item Заявитель получает $d$ (возврат комиссии) $+\ \text{reward}_{\text{pool}} / \mu$
  \item Арбитр/голосующие получают $\text{reward}_{\text{pool}} -
  \text{reward}_{\text{pool}} / \mu$
  \item Остаток страховки: дополнительный штраф $\rightarrow$ фонд DAO
\end{itemize}

\textbf{Оракул прав (спор отклонён):}

\[
d \rightarrow \text{50\% арбитру, 50\% оракулу}
\]

\textbf{Оракул не отвечает (авто-закрытие на 14-й день):}

\begin{itemize}
  \item Заявитель: комиссия возвращена
  \item Оракул: $d$ списывается из страховки
  \item Все ставки и позиции LP: полный возврат
  \item Списанная сумма распределяется пропорционально всем участникам
\end{itemize}

\subsection{Теоретико-игровой анализ}\label{sec:gametheory}

\textbf{Совместимость стимулов оракула.} Ожидаемый выигрыш оракула от честного
разрешения против манипуляции:

\[
\mathbb{E}[\pi_{\text{honest}}] = f_{\text{fixed}} + f_{\text{oracle}} \quad
\text{(за рынок)}
\]

\[
\mathbb{E}[\pi_{\text{dishonest}}] = p_{\text{detect}} \cdot (-I_o \cdot
\theta_{\text{penalty}} - B_{\text{ban}}) + (1 - p_{\text{detect}}) \cdot
g_{\text{manipulation}}
\]

где $p_{\text{detect}}$ --- вероятность спора, $\theta_{\text{penalty}}$ --- доля
списания страховки, $B_{\text{ban}}$ --- приведённая стоимость бана (упущенный
будущий доход), а $g_{\text{manipulation}}$ --- разовый выигрыш от манипуляции.

Честное поведение является равновесием Нэша, когда:

\[
f_{\text{fixed}} + f_{\text{oracle}} > (1 - p_{\text{detect}}) \cdot
g_{\text{manipulation}} - p_{\text{detect}} \cdot (I_o \cdot \theta_{\text{penalty}}
+ B_{\text{ban}})
\]

Протокол обеспечивает это, требуя $I_o \gg g_{\text{manipulation}}$ и поддерживая
высокий $p_{\text{detect}}$ через публичную видимость споров и стимулы участия
комитета.

\textbf{Стимул заявителя.} Рациональный беттор подаёт спор, когда ожидаемое
вознаграждение превышает комиссию спора:

\[
\mathbb{E}[\text{reward}_{\text{dispute}}] = p_{\text{upheld}} \cdot (d + d \cdot
(\mu - 1) / \mu) > d
\]

Это упрощается до $p_{\text{upheld}} > 1/\mu$, задавая естественный порог подачи
спора.

\textbf{Стимул голосования комитета.} В режиме комитета голосующие участвуют, потому
что их стейк в «ленивом» пуле (и, следовательно, вес управления) напрямую
зарабатывает долю пула вознаграждения арбитра. Ожидаемое вознаграждение голосующего
$v$:

\[
\mathbb{E}[\pi_v] = \frac{w_v}{\sum w_j} \cdot \text{voter\_reward\_pool} \cdot
p_{\text{upheld}}
\]

\subsection{Объявление no-contest}

Оракул, неспособный верифицировать исход, может объявить no-contest по сниженной
стоимости ($50\%$ штрафа спора из страховки). Это создаёт трёхуровневый градиент
стимулов:

\begin{center}
\begin{tabular}{lll}
\hline
\textbf{Действие} & \textbf{Стоимость для оракула} & \textbf{Риск бана} \\
\hline
Добровольный no-contest & $0.5 \cdot d$ из страховки & Нет \\
Проигрыш спора & $I_o \cdot \theta_{\text{penalty}} + \text{доп.}$ & Да \\
Пропуск дедлайна & $I_o \cdot \theta_{\text{miss}}$ & Нет \\
\hline
\end{tabular}
\end{center}

Само объявление no-contest можно оспорить, причём арбитр выбирает из трёх исходов
(победил A, победил B или подтвердить no-contest), что предотвращает злоупотребление.

% =====================================================================
\section{Токен VIZ в эксперименте рынка предсказаний}\label{sec:token}

\subsection{Полезность токена}

Токен VIZ выполняет четыре различные функции в экосистеме протокола Onix:

\begin{enumerate}
  \item \textbf{Средство обмена.} Все ставки, депозиты LP, страховка оракулов и
  комиссии споров номинированы в VIZ. Протокол никогда не эмитирует токены --- он
  строго нулевой по сумме на уровне консенсуса.
  \item \textbf{Вес управления.} VIZ, застейканный как вестинг-доли (или внесённый в
  «ленивый» пул), даёт право голоса в спорах комитета DAO и в управлении параметрами
  цепи. Это создаёт прямую полезность владения токеном помимо спекуляции.
  \item \textbf{Залог оракула.} Страховые залоги оракулов номинированы в VIZ. Залог
  должен превышать потенциальную прибыль оракула от манипуляции, создавая спрос на
  накопление токена операторами оракулов.
  \item \textbf{Эскроу спора.} Комиссии споров блокируются в эскроу в VIZ, создавая
  издержку для несерьёзных споров и механизм вознаграждения для обоснованных вызовов.
\end{enumerate}

\subsection{Петля обратной связи спроса на токен}

Протокол создаёт структурный цикл спроса:

\begin{align*}
\text{депозиты LP} &\xrightarrow{\text{блокировка}} \text{снижение оборотного предложения} \\
\text{страховка оракулов} &\xrightarrow{\text{блокировка}} \text{снижение оборотного предложения} \\
\text{активные ставки} &\xrightarrow{\text{блокировка}} \text{снижение оборотного предложения} \\
\text{вес управления} &\xrightarrow{\text{полезность}} \text{спрос на стейкинг}
\end{align*}

Суммарное заблокированное предложение в любой момент:

\[
S_{\text{locked}} = S_{\text{pool}} + \sum_o I_o + \sum_{\text{active bets}} a_i
+ \sum_{\text{disputes}} d_k
\]

где $S_{\text{pool}}$ --- суммарный заблокированный для вывода баланс «ленивого»
пула (как свободная, так \emph{и} размещённая части заблокированы на период
блокировки депозита), $I_o$ --- страховка оракулов, $a_i$ --- ставки активных
позиций, а $d_k$ --- комиссии споров в эскроу. Это заблокированное предложение
снижает доступное оборотное предложение, потенциально создавая повышательное
давление на цену по мере роста использования протокола --- та же ставка, что делает
любой нативный токен протокола, здесь явно привязанная к измеримой полезности.

\subsection{Экспериментальная ценность токена}

Токен VIZ в этом эксперименте служит \textbf{измерительным инструментом}:

\begin{itemize}
  \item \textbf{Цена как сигнал.} Движения цены токена в ответ на события протокола
  (запуски рынков, разрешения с высоким объёмом, исходы споров) дают непрерывную
  рыночную оценку воспринимаемой ценности протокола.
  \item \textbf{Уровень участия в управлении.} Доля держателей токена, участвующих в
  голосованиях по спорам, измеряет жизнеспособность арбитража на основе DAO.
  \item \textbf{Скорость накопления «ленивого» пула.} Темп депозитов в пул измеряет
  розничный спрос на безрисковую доходность LP, напрямую проверяя гипотезу Q3.
\end{itemize}

% =====================================================================
\section{Полный цикл рынка: подробный пример}\label{sec:example}

Проследим полный жизненный цикл бинарного рынка, чтобы проиллюстрировать работу
протокола.

\subsection{Постановка}

\begin{itemize}
  \item \textbf{Рынок:} «Произойдёт ли событие X к дате Y?» (Бинарный: Да/Нет)
  \item \textbf{Оракул:} $o$ со страховкой $I_o = 5000$ VIZ
  \item \textbf{Начальная ликвидность:} $L = 200$ VIZ от создателя
  \item \textbf{Параметры комиссий:} $\theta_{\text{oracle}} = 50$ bp,
  $\theta_{\text{creator}} = 50$ bp, $\theta_{\text{liq}} = 100$ bp
  \item \textbf{«Ленивый» пул:} $B_{\text{free}} = 10{,}000$ VIZ, $\alpha = 2\%$,
  размещает $a_{\text{alloc}} = 200$ VIZ
\end{itemize}

\begin{remark}[Единицы]
Суммы показаны в VIZ для читаемости, но протокол хранит и вычисляет их как
\textbf{целые милли-VIZ (mVIZ)}, где $1\text{ VIZ} = 1000\text{ mVIZ}$, и каждое
деление округляется вниз (ср. замечание к Теореме~\ref{thm:lp}). Строки комиссий
ниже показаны в mVIZ, чтобы $\lfloor\cdot\rfloor$ был точен; эквивалент в VIZ
приведён рядом. Это устраняет кажущееся несоответствие округления, возникающее при
округлении над целыми VIZ (например, $\lfloor 80\times 50/10000\rfloor$ равно $0$ в
VIZ, но $400$ mVIZ $=0.4$ VIZ в реальных единицах).
\end{remark}

\subsection{Инициализация}

\[
R_A = 100, \quad R_B = 100, \quad k = 10{,}000
\]

«Ленивый» пул авто-размещает 200 VIZ как дополнительный LP:

\[
R_A = 200, \quad R_B = 200, \quad k = 40{,}000
\]

\subsection{Фаза ставок}

\begin{itemize}
  \item \textbf{Алиса} ставит $a_1 = 50$ VIZ на «Да» (сторона A):
  \[
  R'_A = 250, \quad R'_B = \lfloor 40000/250 \rfloor = 160, \quad w_1 = 200 -
  160 = 40
  \]
  \item \textbf{Боб} ставит $a_2 = 80$ VIZ на «Нет» (сторона B):
  \[
  R'_B = 240, \quad R'_A = \lfloor 40000/240 \rfloor = 166, \quad w_2 = 250 -
  166 = 84
  \]
\end{itemize}

Подразумеваемая вероятность: $P(\text{Да}) = 166 / 406 \approx 41\%$, $P(\text{Нет}) =
240 / 406 \approx 59\%$.

\subsection{Разрешение}

Оракул объявляет победу \textbf{«Да»}. Алиса --- единственный победитель; Боб теряет
80 VIZ.

\begin{align*}
S_{\text{lose}} &= 80{,}000 \text{ mVIZ} \;(80 \text{ VIZ}) \\
f_{\text{oracle}} &= \lfloor 80{,}000 \times 50 / 10000 \rfloor = 400 \text{ mVIZ} \;(0.4 \text{ VIZ}) \\
f_{\text{creator}} &= \lfloor 80{,}000 \times 50 / 10000 \rfloor = 400 \text{ mVIZ} \;(0.4 \text{ VIZ}) \\
f_{\text{liq}} &= \lfloor 80{,}000 \times 100 / 10000 \rfloor = 800 \text{ mVIZ} \;(0.8 \text{ VIZ}) \\
W &= 80{,}000 - 400 - 400 - 800 = 78{,}400 \text{ mVIZ} \;(78.4 \text{ VIZ})
\end{align*}

\textbf{Выплата Алисе} (единственный победитель, $w_1 / w_1 = 1$):

\[
\pi_1 = 50{,}000 + \lfloor 78{,}400 \times 40/40 \rfloor - \tau_1 = 128{,}400
\text{ mVIZ} - \tau_1 \;(128.4 \text{ VIZ} - \tau_1)
\]

\textbf{Возврат LP:}

\begin{itemize}
  \item LP создателя: $200 \text{ VIZ основного капитала} + \text{взвешенная по времени доля пула комиссий }
  0.8 \text{ VIZ}$
  \item LP «ленивого» пула: $200 \text{ VIZ основного капитала} + \text{взвешенная по времени доля пула комиссий }
  0.8 \text{ VIZ}$
\end{itemize}

\textbf{Проверка} (Теорема~\ref{thm:lp}):

\begin{align*}
\text{Money}_{\text{IN}} &= 200 + 200 + 50 + 80 = 530 \\
\text{Money}_{\text{OUT}} &= 200 + 200 + 50 + 78.4 + 0.4 + 0.4 + 0.8 = 530 \quad
\checkmark
\end{align*}

\subsection{Сценарий спора}

Если Боб оспаривает в течение $\Delta t_{\text{grace}} = 12\text{ч}$, заплатив $d = 10$
VIZ:

\begin{enumerate}
  \item Все выплаты замораживаются.
  \item Оракул должен ответить в течение 12 часов.
  \item В режиме комитета все держатели токена голосуют (взвешенно по стейку,
  публично, пересматриваемо).
  \item Если спор удовлетворён: выплаты пересчитываются с исправленным исходом,
  страховка оракула списывается.
  \item Если спор отклонён: исходные выплаты идут как есть, Боб теряет 10 VIZ
  (делятся 50/50 между оракулом и арбитром).
  \item Если за 14 дней нет разрешения: авто-закрытие с полными возвратами.
\end{enumerate}

\subsection{Сценарий раннего выхода}\label{sec:exitexample}

Предположим, что вместо удержания Алиса отменяет свою ставку на «Да» в середине
рынка после того, как ставка Боба сдвинула кривую в её пользу, так что возврат по
кривой превышает её ставку в 50 VIZ. По механизму отложенного требования
(\Cref{sec:earlyexit}) ей немедленно возвращается основной капитал,
$\text{refund} = \min(\text{curve\_refund}, 50) = 50$ VIZ, а хвост прибыли
$\max(\text{curve\_refund} - 50, 0)$ записывается как отложенное требование на «Да»,
а не выплачивается из кривой. При разрешении, если «Да» побеждает, её требование
оплачивается FIFO из капнутого бакета $\lfloor 0.33 \times S_{\text{lose}} \rfloor$;
если этот бакет уже опустошён более ранними выходящими, она получает урезание; если
«Нет» побеждает, её требование платит ноль. Основной капитал LP не затронут ни в
одной ветке, а победители, \emph{удержавшие} позицию до разрешения, получают не менее
$(1 - 0.33)\,S_{\text{lose}}$ пула проигравших плюс любой неиспользованный бакет ---
конкретно иллюстрируя Теорему~\ref{thm:earlyexit}.

% =====================================================================
\section{Анти-MEV: пакетные и commit-reveal ставки}\label{sec:antimev}

\subsection{MEV на рынках предсказаний}

На рынках предсказаний в непрерывном времени опережающее исполнение (front-running)
и сэндвич-атаки извлекают ценность у бетторов. Когда крупная ставка транслируется в
mempool, атакующий может:

\begin{enumerate}
  \item Опередить: разместить ставку на той же стороне перед крупной ставкой,
  наживаясь на движении цены
  \item Сэндвич: разместить ставки на обеих сторонах вокруг крупной ставки
\end{enumerate}

\subsection{Пакетный расчёт}

Протокол поддерживает опциональные пакетные ставки (бинарные рынки). Ставки,
поданные в пределах эпохи из $E$ блоков, ставятся в очередь и рассчитываются по
\textbf{единой цене} на границе эпохи. Только чистый остаток (разница совокупного
спроса) сдвигает AMM:

\[
\Delta R_A = \sum_{\text{batch}} a_{i,A} - \sum_{\text{batch}} a_{i,B}
\]

Это устраняет преимущество порядка внутри пакета: все ставки пакета получают
идентичное ценообразование независимо от порядка подачи.

\subsection{Commit-reveal}

Для более сильной защиты от MEV бетторы могут использовать двухфазную схему
commit-reveal:

\begin{enumerate}
  \item \textbf{Commit:} Подать $H(\text{bet} \| \text{nonce})$ с эскроу.
  Направление и сумма ставки скрыты.
  \item \textbf{Reveal:} После закрытия эпохи подать $(\text{bet},
  \text{nonce})$ для исполнения по пакетной цене.
\end{enumerate}

Нераскрытые обязательства теряют заданный управлением процент штрафа от эскроу, что
предотвращает спам-обязательства.

% =====================================================================
\section{Реализация}\label{sec:impl}

\subsection{Операции консенсусного уровня в VIZ DLT}

Протокол реализован как первоклассные операции, валидируемые консенсусом, в
распределённом реестре VIZ --- не смарт-контракты, не пользовательские полезные
нагрузки. VIZ DLT обеспечивает время блока ${\sim}3$ секунды, консенсус Delegated
Proof of Stake и именованные аккаунты (в стиле Graphene) без виртуальной машины
общего назначения.

Каждое финансовое действие (создание рынка, размещение ставки, разрешение оракулом,
подача спора, депозит/вывод LP) --- это операция \texttt{pm\_*}, валидируемая каждым
узлом-валидатором. Невалидные операции отклоняются до включения в блок.

\begin{center}
\begin{tabular}{lll}
\hline
\textbf{Слой} & \textbf{Примеры} & \textbf{Валидируется консенсусом} \\
\hline
Операции протокола & \texttt{pm\_create\_market}, \texttt{pm\_place\_bet}, \ldots & Да --- валидирует каждый узел \\
Виртуальные операции & \texttt{pm\_payout}, \texttt{pm\_early\_exit\_claim\_paid}, \ldots & Да --- детерминированно, на момент блока \\
Метаданные & Описания рынков, доказательства споров & Нет --- только на клиенте \\
\hline
\end{tabular}
\end{center}

\subsection{Состояние в цепи}\label{sec:onchainstate}

Всё состояние протокола хранится в индексированных объектах chainbase. Ключевые
объекты:

\begin{itemize}
  \item \texttt{pm\_market\_object}: конфигурация рынка, резервы CPMM, параметры комиссий, состояние
  \item \texttt{pm\_bet\_object}: позиции бетторов с весом и временным штрафом
  \item \texttt{pm\_liquidity\_object}: позиции LP с весом по времени
  \item \texttt{pm\_lazy\_pool\_object}: состояние синглтон-пула (балансы, доли, накопитель вознаграждения)
  \item \texttt{pm\_oracle\_object}: регистрация оракула, страховка, счётчики репутации
  \item \texttt{pm\_dispute\_object}: состояние и разрешение спора
  \item \texttt{pm\_deferred\_claim\_object}: зависящие от исхода требования раннего выхода (FIFO по времени выхода)
\end{itemize}

Метрики репутации вычисляются при чтении (не хранятся), что обеспечивает
согласованность без дополнительных операций записи.

\subsection{Параметры управления}\label{sec:governance}

Все экономические параметры (комиссии, штрафы, требования к страховке, окна споров,
настройки «ленивого» пула и кап вознаграждения раннего выхода
\texttt{pm\_early\_exit\_reward\_cap\_percent}) выбираются медианным голосованием
делегатов. Каждый избранный валидатор публикует предпочтительные значения; сеть
вычисляет медиану. Параметры меняются без хардфорков и развёртываний. Аварийные
выключатели (kill-switch) позволяют управлению отключать подсистемы (плечо,
commit-reveal) медианным голосованием.

\subsection{Масштабируемость}\label{sec:scalability}

Поскольку каждая операция валидируется консенсусом, модель издержек протокола важна
при масштабе (тысячи одновременных рынков и ставок). Дизайн удерживает работу на блок
и на операцию ограниченной:

\begin{itemize}
  \item \textbf{Операции пользователя за $O(1)$.} Размещение ставки, отмена и
  депозит/вывод LP затрагивают фиксированное число индексированных объектов (рынок,
  позицию беттора и --- для пула --- синглтон-накопитель). Ни одна не итерирует по
  всем участникам. Распределение вознаграждений использует накопитель MasterChef
  $\rho$ (\Cref{sec:lazyacct}), так что пул с $n$ вкладчиками рассчитывает
  вознаграждения за $O(1)$ на участника, а не за $O(n)$ на разрешение.
  \item \textbf{Ограниченная отложенная работа на блок.} Расчёт --- единственный шаг
  с разветвлением (рынок с $m$ выигравшими ставками порождает $m$ виртуальных
  операций \texttt{pm\_payout}; распределение раннего выхода добавляет не более одной
  \texttt{pm\_early\_exit\_claim\_paid} на каждое сохранившееся отложенное требование,
  само по себе капнутое на рынок). Чтобы одно крупное разрешение не раздувало блок,
  обработка выплат и cron ограничена медианным параметром
  \texttt{pm\_processing\_cap\_per\_block}: за блок обрабатывается не более
  фиксированного числа выплат/cron-элементов, а остаток переносится на последующие
  блоки. Курсор round-robin по рынкам гарантирует, что новые рынки не голодают из-за
  «старожилов» с низким id. Поэтому работа на блок в худшем случае --- $O(\text{cap})$,
  независимо от того, сколько рынков разрешается в одном интервале.
  \item \textbf{Нет усиления записи от репутации.} 14 метрик оракула --- счётчики,
  обновляемые только при собственных действиях оракула; составной балл надёжности
  вычисляется \textbf{при чтении}, а не пишется каждый блок (\Cref{sec:onchainstate}).
  Метаданные обнаружения рынков строятся неконсенсусным плагином и никогда не входят
  в валидацию блока.
  \item \textbf{Рост состояния.} Состояние линейно по числу живых объектов (рынки,
  открытые позиции, активные споры, нерассчитанные отложенные требования). Разрешённые
  рынки и оплаченные позиции терминальны и обрезаемы клиентами; консенсус хранит лишь
  то, что требуют открытые обязательства.
\end{itemize}

Связывающее ограничение при очень большом числе рынков --- это пространство блока для
разветвления расчётов, которое \texttt{pm\_processing\_cap\_per\_block} превращает из
всплеска задержки в ограниченную, амортизированную пропускную способность, а не в
стоимость, останавливающую консенсус. Количественная стресс-симуляция по многим
одновременным рынкам разного объёма и волатильности входит в экспериментальную
программу (\Cref{sec:hypotheses}, H4).

% =====================================================================
\section{Обсуждение}\label{sec:discussion}

\subsection{Сравнение с существующими подходами}

\begin{center}
\begin{tabular}{lllll}
\hline
\textbf{Измерение} & \textbf{Onix} & \textbf{Polymarket} & \textbf{Станд. LMSR} & \textbf{AMM Uniswap} \\
\hline
Риск LP & \textbf{Нулевой} (структурно) & Инвентарный риск & До $b \ln N$ & Непостоянные потери \\
Знания LP & Низкие (депозит) & Высокие (управлять заявками) & Средние & Средние-высокие \\
Непрерывность цены & Непрерывная & Дискретная (книга) & Непрерывная & Непрерывная \\
Извлечение комиссий & Только проигравшие, при разрешении & Спред bid-ask & Спред & Каждая сделка \\
Модель оракула & Залог + спор DAO & Оракул UMA & Оператор & --- \\
Согласованность цен & Матем. инвариант & Зависит от арбитража & Матем. инвариант & Матем. инвариант \\
Split/merge CTF & Нет & Да & Нет & Нет \\
\hline
\end{tabular}
\end{center}

\subsection{Ограничения и честные компромиссы}\label{sec:limitations}

\begin{enumerate}
  \item \textbf{Доходность LP зависит от объёма, а не от глубины.} Рынок с высокой
  субсидией и низким объёмом зарабатывает столько же абсолютных комиссий, сколько
  рынок с низкой субсидией и равным объёмом. Субсидия даёт глубину (меньше
  проскальзывания), но не доходность.
  \item \textbf{Прибыль LP не гарантирована.} Если рынок разрешается с нулевыми
  проигравшими ставками, комиссии не образуются. Основной капитал LP возвращается, но
  доходность может быть нулевой.
  \item \textbf{Тотализаторные токены --- не инструменты с фиксированной стоимостью.}
  В отличие от стандартного LMSR, где 1 выигрышный токен $=$ 1 единице валюты, токены
  Onix --- пропорциональные требования на пул проигравших. Если все бетторы выбрали
  победителя, все выходят «в ноль».
  \item \textbf{Компромиссы управления DPoS.} Модель параметров, выбираемых
  делегатами, имеет известные риски концентрации, присущие DPoS (общие с EOS, Hive,
  Tron). Механизм веса управления «ленивого» пула частично смягчает это, расширяя
  эффективный электорат.
  \item \textbf{Зависимость от ликвидности токена.} Экономические гарантии (страховые
  залоги, комиссии споров) масштабируются с рыночной стоимостью токена. Протокол
  предполагает, что полезность движет спросом, --- стандартное допущение для нативных
  токенов протоколов.
  \item \textbf{Плечо добавляет ограниченный кредитный риск пулу.} Опциональная
  подсистема плеча (\Cref{sec:leverage}) --- единственный компонент, \emph{не}
  покрытый Теоремой~\ref{thm:lp}. Будучи включённой, пул действует как кредитор и
  может понести недостачу на пути отмены той же стороны, ограниченную залогом заёмщика
  и изолированную в управляемой доле фонда. Мы заявляем полный возврат ссуды на путях
  каскада и расчёта как \emph{проектное свойство, проверенное в симуляции}, а не как
  доказанную теорему; замкнутое доказательство возврата при произвольных траекториях
  резервов --- открытая задача. Подсистема выключена по умолчанию и отключаема
  медианным голосованием, поэтому безусловная гарантия LP всегда может быть
  восстановлена.
  \item \textbf{Ранний выход --- это ограниченный дисконт, и кап --- вопрос политики.}
  Механизм отложенного требования (\Cref{sec:earlyexit}) даёт \emph{жёсткую},
  безусловную гарантию LP (Теорема~\ref{thm:earlyexit}, через ограничение по
  headroom), но ценой того, что ранний выход становится капнутой, упорядоченной по
  FIFO, зависящей от исхода выплатой, а не обналичиванием по кривой (mark-to-curve).
  Кап $c$ --- это выбор политики управления: слишком низкий кап делает ранний выход
  непривлекательным и истончает вторичную ликвидность, слишком высокий кап размывает
  награду держателям --- а в углу с высокими комиссиями ограничение расчёта далее
  сжимает эффективный бакет, чтобы сохранить платёжеспособность, углубляя урезание
  раннего выхода. Нахождение значения, которое максимизирует вторичную ликвидность, не
  отбивая охоту удерживать, --- эмпирический вопрос (\Cref{sec:hypotheses}, H8), а не
  тот, что решает теория.
  \item \textbf{Публичное голосование по спорам имеет компромисс тайминга.} Открытые,
  пересматриваемые бюллетени комитета (\Cref{sec:disputemodes}) допускают преимущество
  последнего хода и выбраны намеренно вместо commit-reveal ради аудируемости. Это
  ценностное суждение (прозрачность важнее тайны), а не доказательство, что открытое
  голосование оптимально против манипуляций; развёртывания, ставящие тайну в
  приоритет, должны вместо этого использовать режим аккаунта.
\end{enumerate}

\subsection{Экспериментальные гипотезы}\label{sec:hypotheses}

Развёрнутый протокол спроектирован для проверки следующих гипотез:

\textbf{H1 (Маховик ликвидности).} \emph{Безрисковое предоставление ликвидности
привлекает розничный капитал, достаточный для создания рыночной глубины, которая
ощутимо снижает проскальзывание относительно сопоставимых платформ.} Измеримо:
суммарные депозиты «ленивого» пула, средняя глубина рынка (отношения резервов),
проскальзывание на единицу размера ставки.

\textbf{H2 (Агрегирование информации).} \emph{Ценообразование AMM/LMSR с
тотализаторным расчётом даёт вероятностные оценки точности, сопоставимой с рынками
предсказаний на основе CLOB при эквивалентных информационных условиях.} Измеримо:
баллы Бриера, кривые калибровки, сравнение с внешними бенчмарками.

\textbf{H3 (Надёжность оракула под управлением DAO).} \emph{Оракулы с залогом при
разрешении споров в режиме комитета достигают точности разрешения, сопоставимой с
централизованными оракульными сервисами, при уровне споров ниже устойчивого порога.}
Измеримо: уровень споров оракула, уровень проигрыша споров, среднее время разрешения,
распределение баллов надёжности.

\textbf{H4 (Устойчивость «ленивого» пула).} \emph{Механизмы градуированного отзыва и
штрафных меток предотвращают эксплуатацию «ленивого» пула оракулами, поддерживая
положительную чистую доходность вкладчиков пула по разнообразным портфелям рынков.}
Измеримо: NAV пула во времени, частота отзыва, распределение штрафных меток, чистая
доходность на долю.

\textbf{H5 (Корреляция спроса на токен).} \emph{Использование протокола (объём, число
рынков, депозиты LP) положительно коррелирует со спросом на стейкинг токена и уровнем
участия в управлении.} Измеримо: уровень стейкинга токена, участие в голосованиях по
спорам, темп депозитов «ленивого» пула, корреляционный анализ.

\textbf{H6 (Эффективность анти-MEV).} \emph{Пакетный расчёт и механизмы commit-reveal
снижают измеримое извлечение MEV по сравнению с непрерывными мгновенными ставками.}
Измеримо: асимметрия влияния на цену (пакет против мгновенной), частота сэндвич-атак,
качество исполнения для беттора.

\textbf{H7 (Платёжеспособность плеча).} \emph{При живых траекториях ликвидация на
основе снимка возвращает ссуду на путях каскада и расчёта, ограничивая реализованную
недостачу пула путём отмены той же стороны и в пределах залога заёмщика, так что
процент по плечу чисто-приращивает доходность пула.} (Это эмпирически проверяет
проектное свойство из \Cref{sec:leverage} вместо замкнутого доказательства возврата.)
Измеримо: фактический коэффициент возврата ссуды на событие ликвидации
($\text{recovered}/\lambda m$), частота и величина недостач отмены той же стороны
относительно внесённого залога, доля процента по плечу в общей доходности пула, NAV
пула с включённой и выключенной подсистемой плеча.

\textbf{H8 (Дисконт ликвидности при раннем выходе).} \emph{Механизм отложенного
требования (\Cref{sec:earlyexit}) обеспечивает полезную вторичную ликвидность ---
бетторы и позиции с плечом действительно выходят рано, --- при этом капнутый бакет
FIFO держит выплату раннего выхода строго ниже выплаты за удержание, так что
удерживаемые победители не размываются и ни один ранний выход никогда не черпает из
основного капитала LP.} Измеримо: объём и частота раннего выхода, распределение
отношений \texttt{paid}/\texttt{claimed} (урезание) из виртуальной операции
\texttt{pm\_early\_exit\_claim\_paid}, реализованная доходность раннего выхода против
доходности удержания на сопоставленных позициях и эмпирический уровень
$\text{uncovered}$ (ожидание: тождественно ноль) как живая проверка
Теоремы~\ref{thm:earlyexit}.

% =====================================================================
\section{Заключение}\label{sec:conclusion}

Мы представили протокол Onix --- архитектуру рынка предсказаний, которая достигает
структурной гарантии основного капитала LP, разделяя функцию ценообразования (CPMM
для бинарного, LMSR для многоисходного) и функцию расчёта (тотализатор, проигравшие
финансируют выигравших). Мы доказали, что эта гарантия выполняется для обоих типов
рынков независимо от назначения весов, параметров комиссий или распределения ставок
(Теорема~\ref{thm:lp}), и что она \emph{переживает присутствие кривой, которая также
ценообразует ранние выходы}, --- механизм отложенного права требования, зависящего от
исхода (Теорема~\ref{thm:earlyexit}), позволяет бетторам и позициям с плечом уходить
рано, ни разу не реализуя торговую P\&L против глубины LP, без непокрытой недостачи и
без эмиссии.

Механизм «ленивого» пула обеспечивает пассивное участие розничных LP с автоматическим
размещением капитала и градуированной защитой от издержек упущенных возможностей.
Двухрежимная система споров --- сочетающая оракулов с залогом со взвешенным по стейку
голосованием комитета DAO --- даёт жизнеспособную альтернативу централизованному
арбитражу, сохраняя совместимость стимулов оракула.

Протокол реализован как операции консенсусного уровня в распределённом реестре VIZ и
спроектирован как эксперимент, проверяющий восемь явных гипотез о запуске
ликвидности, точности агрегирования информации, жизнеспособности управления DAO,
платёжеспособности плеча, ликвидности раннего выхода и динамике спроса на токен.
Опциональная подсистема плеча --- единственный компонент, который меняет безусловную
гарантию LP на ограниченный, выключенный по умолчанию кредитный риск пула; всё
остальное сохраняет структурные гарантии Теорем~\ref{thm:lp}
и~\ref{thm:earlyexit}.

Математические свойства проверяемы. Экономические гипотезы будут проверены участием
рынка.

% =====================================================================
\begin{thebibliography}{99}

\bibitem{hanson2003} Hanson, R. (2003). Combinatorial Information Market Design.
\emph{Information Systems Frontiers}, 5(1), 107--119.
\url{https://doi.org/10.1023/A:1022058209073}

\bibitem{uniswapv2} Adams, H., Zinsmeister, N., \& Robinson, D. (2020).
\emph{Uniswap v2 Core.} Uniswap Labs. \url{https://uniswap.org/whitepaper.pdf}

\bibitem{uniswapv3} Adams, H., Zinsmeister, N., Salem, M., Keefer, R., \&
Robinson, D. (2021). \emph{Uniswap v3 Core.} Uniswap Labs.
\url{https://uniswap.org/whitepaper-v3.pdf}

\bibitem{berg2008} Berg, J., Forsythe, R., Nelson, F., \& Rietz, T. (2008).
Results from a Dozen Years of Election Futures Markets Research. In
\emph{Handbook of Experimental Economics Results} (Vol. 1, pp. 742--751).
Elsevier. \url{https://doi.org/10.1016/S1574-0722(07)00080-7}

\bibitem{cowgill2009} Cowgill, B., Wolfers, J., \& Zitzewitz, E. (2009). Using
Prediction Markets to Track Information Flows: Evidence from Google. In
\emph{Auctions, Market Mechanisms and Their Applications (AMMA 2009)}, LNICST
Vol. 14. Springer. \url{https://doi.org/10.1007/978-3-642-03821-1_2}

\bibitem{levitt2004} Levitt, S. D. (2004). Why Are Gambling Markets Organised So
Differently from Financial Markets? \emph{The Economic Journal}, 114(495),
223--246. \url{https://doi.org/10.1111/j.1468-0297.2004.00207.x}

\bibitem{gnosisctf} Gnosis. \emph{Conditional Tokens Framework (CTF)
Documentation.} \url{https://docs.gnosis.io/conditionaltokens/}

\bibitem{uma} UMA Protocol. \emph{Optimistic Oracle Documentation.}
\url{https://docs.uma.xyz/}

\bibitem{compound} Leshner, R., \& Hayes, G. (2019). \emph{Compound: The Money
Market Protocol.} Compound Labs.
\url{https://compound.finance/documents/Compound.Whitepaper.pdf}

\bibitem{masterchef} SushiSwap. (2020). \emph{MasterChef Contract.}
\url{https://github.com/sushiswap/sushiswap}

\bibitem{viz} Piskunov, A. \emph{VIZ: Distributed Ledger Technical Description,
Fair DPoS, and Governance.} VIZ-Blockchain.
\url{https://viz-blockchain.github.io/viz-cpp-node/}

\bibitem{arrow2008} Arrow, K. J., Forsythe, R., Gorham, M., Hahn, R., Hanson,
R., Ledyard, J. O., et al. (2008). The Promise of Prediction Markets.
\emph{Science}, 320(5878), 877--878. \url{https://doi.org/10.1126/science.1157679}

\bibitem{wolfers2004} Wolfers, J., \& Zitzewitz, E. (2004). Prediction Markets in
Theory and Practice. \emph{NBER Working Paper} No. 10248 (also \emph{Journal of
Economic Perspectives}, 18(2), 107--126). \url{https://doi.org/10.3386/w10248}

\bibitem{othman2013} Othman, A., Pennock, D. M., Reeves, D. M., \& Sandholm, T.
(2013). A Practical Liquidity-Sensitive Automated Market Maker. \emph{ACM
Transactions on Economics and Computation}, 1(3), Article 14.
\url{https://doi.org/10.1145/2509413.2509414}

\bibitem{abernethy2011} Abernethy, J., Chen, Y., \& Vaughan, J. W. (2011). An
Optimization-Based Framework for Automated Market-Making. In \emph{Proceedings of
the 12th ACM Conference on Electronic Commerce (EC '11)} (pp. 297--306).
\url{https://doi.org/10.1145/1993574.1993621}

\end{thebibliography}

\end{document}
